\subsection{Answers to the Research Questions}
\label{subsec:beantwortung-fragestellungen}

\paragraph{Research Question 1: Requirements Addressed.}
The template addresses the functional and technical, quality, and operational
requirements through a baseline that can be generated, clearly defined
components, profiles, shared workflows, centralized maintenance, and versioned
configuration. Separate runtime and verification paths support this;
domain logic, authentication, authorization, and production operation remain
product-specific.

\paragraph{Research Question 2: Variants for Web, Cloud, and Desktop.}
Five presets combine the shared frontend with backend, cloud, and Tauri
components; PostgreSQL is added only when backend capability is present.
Generated projects predominantly contain the paths they require.
\texttt{desktop-cloud} and \texttt{full-platform} are scaffolded identically
at the technical level and differentiated semantically.

\paragraph{Research Question 3: Reuse and Standardization.}
Centralized maintenance of implementations, profiles, the generator,
configuration, tests, and documentation promotes reuse and consistent initial
states across projects. The benefit applies to the baseline and future
generations; existing product projects receive no automatic updates.

\paragraph{Research Question 4: Reduction of Recurring Effort.}
Automated or standardized profile selection, generation, diagnostics,
installation, startup, testing, and builds provide plausible potential for
reducing recurring effort. Time, effort, and comparative data are lacking,
however; a productivity improvement has not been demonstrated when central
maintenance and multiple toolchains are taken into account.

\paragraph{Research Question 5: Suitability for Project Types and Customer Requirements.}
The template is suitable as a foundation for related web, desktop, and
combined business applications whose platform, cloud, and persistence needs
can be represented. Highly native applications, games, and specialized
real-time systems lie outside the primary target domain. The qualitative
classification is not a universal assurance of suitability.

\paragraph{Research Question 6: Limitations, Maintenance Effort, Risks, and External Dependencies.}
Limitations include the growing breadth of regression testing, manual contract
synchronization, insufficiently strict separation between optional and
selectable capabilities, the absence of product updates, and the effort of
maintaining multiple toolchains.
The tooling, desktop-profile, and Windows CI paths are currently failing;
Release Validation, Compose startup, and Branch Protection are still missing.
Signing, notarization, the updater, production deployment, and the security
model remain product tasks.
